\documentclass{article}
\usepackage{mathnotes}

\title{Statistics}
\description{Statistics, estimators, and sufficiency.}

\begin{document}

\section{Statistics}

\begin{definition}[Statistical Model]
TODO: define this.
\end{definition}

\begin{definition}[Population Parameter]
TODO: define this.
\end{definition}

\begin{definition}[Statistic]
TODO: define this.
\end{definition}

\begin{definition}[Estimator]
TODO: define this.
\end{definition}

\begin{definition}[Estimate]
TODO: define this.
\end{definition}

\begin{definition}[Bias]
TODO: define this.
\end{definition}

\begin{definition}[Unbiased Estimator]
TODO: define this.
\end{definition}

\subsection{Sufficient Statistics}

\begin{definition}[Sufficient Statistic]
	A \@{function} $T(X)$ is said to be a \textbf{sufficient statistic} relative to the family $\{f_\theta(x)\}$ if $X$ is \@{conditionally independent} of $\theta$ given $T(X)$ for any \@{probability distribution} on $\theta,$ that is, $\theta \to T(X) \to X$ forms a \@{Markov chain}.
\end{definition}

\begin{definition}[Minimal Sufficient Statistic]
	A \@{statistic} $T(X)$ is a \@{minimal sufficient statistic} relative to $\{f_\theta(x)\}$ if it is a \@{function} of every other \@{sufficient statistic} $U.$ 
\end{definition}
\begin{note}
	In terms of the \@{data-processing inequality}, this implies that

	\[ \theta \to T(X) \to U(X) \to X. \]
\end{note}

\end{document}
